\documentclass[11pt]{article}
\usepackage[margin=1.05in,headheight=14pt]{geometry}
\usepackage[T1]{fontenc}
\usepackage{lmodern,microtype,amsmath,amssymb,amsthm,booktabs}
\usepackage{xcolor}
\definecolor{linknavy}{RGB}{25,55,125}
\usepackage[colorlinks=true,linkcolor=linknavy,citecolor=linknavy,urlcolor=linknavy]{hyperref}
\usepackage{orcidlink,fancyhdr}
\newtheorem{theorem}{Theorem}
\newtheorem{lemma}[theorem]{Lemma}
\theoremstyle{remark}
\newtheorem{remark}[theorem]{Remark}
\newcommand{\e}{\mathrm{e}}
\newcommand{\ntriples}{47}
\newcommand{\nconfirmed}{47}
\newcommand{\ngraphs}{6{,}224{,}502{,}963}


% Page-1 header block (conventions/manuscript-header-standard.md). The four document
% macros are set by the publishing tool when the version DOI is reserved.
\newcommand{\doctype}{Preprint}
\newcommand{\docversion}{v0.02}
\newcommand{\docdate}{8 November 2026}
\newcommand{\versiondoi}{10.5281/zenodo.24059677}
\newcommand{\conceptdoi}{10.5281/zenodo.24044537}
\newcommand{\docrepo}{github.com/fsantibanezleal/CAOS\_RESEARCH}
\newcommand{\headerblock}{%
\begin{center}
\fbox{\parbox{0.94\textwidth}{\small\raggedright
\textbf{\doctype} \quad$\cdot$\quad \docversion, \docdate \quad$\cdot$\quad Licence CC BY 4.0\\[2pt]
CAOS Research open-problems programme \quad$\cdot$\quad problem \texttt{bougard-joret}\\[2pt]
Version DOI \href{https://doi.org/\versiondoi}{\texttt{\versiondoi}}\\[2pt]
Concept DOI (always latest) \href{https://doi.org/\conceptdoi}{\texttt{\conceptdoi}}\\[2pt]
Not peer reviewed. Source code, data and computational records: \href{https://\docrepo}{\texttt{\docrepo}}
}}
\end{center}
\vspace{4pt}}

\pagestyle{fancy}
\fancyhf{}
\fancyhead[L]{\small CAOS Research \doctype}
\fancyhead[R]{\small Second regime\textperiodcentered\ \docversion}
\fancyfoot[C]{\thepage}
\fancypagestyle{plain}{\fancyhf{}\fancyfoot[C]{\thepage}\renewcommand{\headrulewidth}{0pt}}
\title{Exact values of the Bougard--Joret function\\in its open second regime up to order 23}
\author{Felipe Santib\'a\~nez-Leal\,\orcidlink{0000-0002-0150-3246}\\
\small CAOS Research program\\
\small\texttt{github.com/fsantibanezleal/CAOS\_RESEARCH}}
\date{\docdate}
\begin{document}
\maketitle
\headerblock
\begin{abstract}
Let $f(n,\alpha,k)$ be the minimum number of edges of a $k$-connected graph of order $n$ with independence number
$\alpha$. Bougard and Joret conjectured $f(n,\alpha,k)=t(n,\alpha)+\lceil k\alpha/2\rceil$ when $n>k\alpha$,
where $t(n,\alpha)$ is the size of the balanced union of $\alpha$ cliques. Das and Gupta disproved the other line
of the conjecture and left this second regime open. Bougard and Joret proved the second line only when
$\lfloor n/\alpha\rfloor\ge\lceil k\alpha/2\rceil-\alpha+2$ (for $\alpha\ge3$). We compute $f$ exactly at every
triple outside that range with $\alpha\ge3$, $k\ge3$ and $n\le23$: \ntriples\ triples. At every one of them the
value equals the conjectured one, by complete enumeration of every candidate graph below it
(\ngraphs\ candidate graphs in all). The enumeration uses nauty with two pruning rules that are proved sound, and an exact filter checked
against an independent implementation. The data support the second line of the conjecture in its smallest open
cases; they do not prove it in general.
\end{abstract}

\section{Problem and prior work}

All graphs are finite and simple. Write $\e(G)$ for the number of edges, $\alpha(G)$ for the independence number
and $\kappa(G)$ for the vertex connectivity. For $n\ge2\alpha$, $n\ge\alpha+k$, $\alpha\ge2$ and $k\ge3$, Bougard
and Joret conjectured \cite[Conjecture 1]{BJ}
\[
f(n,\alpha,k)=\begin{cases}\lceil nk/2\rceil,&n\le k\alpha,\\ t(n,\alpha)+\lceil k\alpha/2\rceil,&n>k\alpha.\end{cases}
\]
The upper bound in the second regime holds by their construction: the balanced union of $\alpha$ cliques
$T(n,\alpha)$ becomes $k$-connected after adding $\lceil k\alpha/2\rceil$ edges without raising the independence
number \cite[Section 6]{BJ}. Das and Gupta determined the boundary $n=\alpha+k$, which disproves the first line
(the smallest counterexample is $f(7,3,4)=15$), and left the second regime open \cite{DG}. The first interior shell
$n=\alpha+k+1$ \cite{Shell} and the value $f(2d+3,2,d+2)=d^2+4d+2$ \cite{Next} lie in the first regime.

\paragraph{What is proved in the second regime.} Bougard and Joret's Proposition 8 states that a 2-connected graph
with $n\ge3\alpha\ge9$ and at most $t(n,\alpha)+\lfloor n/\alpha\rfloor+\alpha-3$ edges has a vertex partition into
$\alpha$ cliques \cite{BJ}. A $k$-connected graph with such a partition has at least
$t(n,\alpha)+\lceil k\alpha/2\rceil$ edges (each clique sends at least $k$ edges out of it). Hence the second line
holds whenever $\lceil k\alpha/2\rceil-1\le\lfloor n/\alpha\rfloor+\alpha-3$, that is
\begin{equation}\label{eq:proved}
\lfloor n/\alpha\rfloor\ \ge\ \lceil k\alpha/2\rceil-\alpha+2\qquad(\alpha\ge3),
\end{equation}
and also for $\alpha=2$, $n\ge2k$ \cite{BJ}. The \emph{open window} is the set of second-regime triples with
$\alpha\ge3$ violating \eqref{eq:proved}. Its smallest members are $(10,3,3)$ and $(11,3,3)$.

\section{Method}

For a triple $(n,\alpha,k)$ with conjectured value $c=t(n,\alpha)+\lceil k\alpha/2\rceil$ we decide, for every
edge count $E<c$ that can occur, whether some graph on $n$ vertices with $E$ edges has independence number
exactly $\alpha$ and connectivity at least $k$. If none does, $f(n,\alpha,k)\ge c$, and with the construction,
$f(n,\alpha,k)=c$. Two classical facts bound the edge counts: minimum degree at least $k$ gives
$E\ge\lceil nk/2\rceil$, and Tur\'an's theorem applied to the complement gives $E\ge t(n,\alpha)$ for independence
number at most $\alpha$ \cite{Turan}. So only $\lceil k\alpha/2\rceil$ edge counts per triple need to be searched.

Every graph with $n$ vertices, $E$ edges and minimum degree at least $k$ is generated up to isomorphism by
\texttt{geng} from nauty~2.8.9 \cite{nauty}, split into disjoint parts. Each graph is tested by an exact filter:
the independence number by branch and bound on bit sets, and connectivity at least $k$ by deleting every vertex set
of size less than $k$. Two refinements were added during the work; each is proved sound and was gated before use.

\begin{lemma}[early stop]\label{lem:early}
Rejecting a graph as soon as the search finds an independent set of size $\alpha+1$ leaves the decision
``$\alpha(G)=\alpha$'' unchanged.
\end{lemma}
\begin{proof}
Such a graph has $\alpha(G)\ge\alpha+1$ and is rejected by the full computation as well; otherwise the search
runs to the end and computes $\alpha(G)$ exactly.
\end{proof}

\begin{lemma}[independence prune]\label{lem:prune}
During generation, a partial graph that already contains an independent set of size $\alpha+1$ can be discarded
with all its extensions.
\end{lemma}
\begin{proof}
\texttt{geng} extends a graph by adding vertices only, so the partial graph is an induced subgraph of each of its
completions. The independence number of an induced subgraph never exceeds that of the graph, so every completion
has independence number at least $\alpha+1$ and is rejected.
\end{proof}

With Lemma~\ref{lem:prune}, the $(17,5,3)$ search at $E=27$ dropped from $1{,}040$ seconds
($35{,}329{,}226$ candidates) to under a second ($482$ candidates), with the same answer.

\section{Validation}

\begin{itemize}
\item The filter agrees with an independent implementation (networkx) on all $12{,}346$ graphs on 8 vertices, for
  every $\alpha\le6$ and $k\le5$, before and after Lemma~\ref{lem:early}.
\item Known values are reproduced: $f(7,3,4)=15$ \cite{DG}, $f(6,2,3)=9$, $f(6,2,4)=12$, and the second-regime
  value $f(12,3,3)=23$ inside the proved range \eqref{eq:proved}.
\item The generator with Lemma~\ref{lem:prune} reproduces the numbers of qualifying graphs found without it at
  twelve known cases, from $(7,3,4)$ to $(17,5,3)$.
\item Connected-graph counts of \texttt{geng} agree with OEIS A001349 where checked, and the 4,207 graphs at
  $(16,5,3)$, $E=24$ are the cubic graphs on 16 vertices (OEIS A005638).
\end{itemize}

\section{Results}

\begin{theorem}[MV]
At each of the \ntriples\ triples of the open window with $n\le23$ (Table~\ref{tab:values}), no graph with fewer
edges than the conjectured value has independence number $\alpha$ and connectivity at least $k$; hence
$f(n,\alpha,k)=t(n,\alpha)+\lceil k\alpha/2\rceil$ at every one of them.
\end{theorem}

\begin{table}[p]
\centering\footnotesize
\begin{tabular}{lrrrrr}
\toprule
$(n,\alpha,k)$ & $t(n,\alpha)$ & $f$ & edges searched & graphs below & extremal graphs \\
\midrule
$(10,3,3)$ & 12 & 17 & 15--16 & 383 & 15 \\
$(11,3,3)$ & 15 & 20 & 17--19 & 55{,}812 & 6 \\
$(13,3,4)$ & 22 & 28 & 26--27 & 1{,}095{,}032 & 1 \\
$(13,4,3)$ & 15 & 21 & 20--20 & 1{,}999 & 1 \\
$(14,3,4)$ & 26 & 32 & 28--31 & 5{,}268{,}961{,}461 & construction \\
$(14,4,3)$ & 18 & 24 & 21--23 & 2{,}147{,}042 & 1 \\
$(15,4,3)$ & 21 & 27 & 23--26 & 904{,}211{,}641 & 1 \\
$(16,3,5)$ & 35 & 43 & 40--42 & 134 & 3 \\
$(16,5,3)$ & 18 & 26 & 24--25 & 932{,}175 & 18 \\
$(17,3,5)$ & 40 & 48 & 43--47 & 2{,}120 & 3 \\
$(17,4,4)$ & 28 & 36 & 34--35 & 99 & 2 \\
$(17,5,3)$ & 21 & 29 & 26--28 & 35{,}627{,}553 & 34 \\
$(18,3,5)$ & 45 & 53 & 45--52 & 6{,}568 & 2 \\
$(18,4,4)$ & 32 & 40 & 36--39 & 4{,}086 & 4 \\
$(18,5,3)$ & 24 & 32 & 27--31 & 71{,}021 & 35 \\
$(19,3,5)$ & 51 & 59 & 51--58 & 17{,}920 & 4 \\
$(19,3,6)$ & 51 & 60 & 57--59 & 83 & 1 \\
$(19,4,4)$ & 36 & 44 & 38--43 & 24{,}673 & 2 \\
$(19,5,3)$ & 27 & 35 & 29--34 & 149{,}597 & 18 \\
$(19,6,3)$ & 21 & 30 & 29--29 & 295 & 2 \\
$(20,3,5)$ & 57 & 65 & 57--64 & 21{,}372 & 4 \\
$(20,3,6)$ & 57 & 66 & 60--65 & 4{,}046 & 1 \\
$(20,4,4)$ & 40 & 48 & 40--47 & 41{,}808 & 2 \\
$(20,5,3)$ & 30 & 38 & 30--37 & 132{,}654 & 5 \\
$(20,6,3)$ & 24 & 33 & 30--32 & 23{,}316 & 5 \\
$(21,3,6)$ & 63 & 72 & 63--71 & 29{,}091 & 1 \\
$(21,4,4)$ & 45 & 53 & 45--52 & 141{,}850 & 2 \\
$(21,4,5)$ & 45 & 55 & 53--54 & 420 & 1 \\
$(21,5,3)$ & 34 & 42 & 34--41 & 471{,}911 & 17 \\
$(21,5,4)$ & 34 & 44 & 42--43 & 1{,}298 & 7 \\
$(21,6,3)$ & 27 & 36 & 32--35 & 303{,}615 & 5 \\
$(22,3,6)$ & 70 & 79 & 70--78 & 83{,}931 & 1 \\
$(22,3,7)$ & 70 & 81 & 77--80 & 782 & construction \\
$(22,4,4)$ & 50 & 58 & 50--57 & 222{,}379 & 4 \\
$(22,4,5)$ & 50 & 60 & 55--59 & 47{,}945 & construction \\
$(22,5,3)$ & 38 & 46 & 38--45 & 874{,}494 & 31 \\
$(22,5,4)$ & 38 & 48 & 44--47 & 101{,}779 & construction \\
$(22,6,3)$ & 30 & 39 & 33--38 & 1{,}365{,}978 & 5 \\
$(22,7,3)$ & 24 & 35 & 33--34 & 20{,}112 & construction \\
$(23,3,6)$ & 77 & 86 & 77--85 & 100{,}793 & construction \\
$(23,3,7)$ & 77 & 88 & 81--87 & 62{,}230 & construction \\
$(23,4,4)$ & 55 & 63 & 55--62 & 190{,}804 & construction \\
$(23,4,5)$ & 55 & 65 & 58--64 & 610{,}557 & construction \\
$(23,5,3)$ & 42 & 50 & 42--49 & 954{,}295 & construction \\
$(23,5,4)$ & 42 & 52 & 46--51 & 1{,}410{,}017 & construction \\
$(23,6,3)$ & 33 & 42 & 35--41 & 2{,}519{,}045 & construction \\
$(23,7,3)$ & 27 & 38 & 35--37 & 1{,}456{,}747 & construction \\
\bottomrule
\end{tabular}

\caption{The open window up to order 23. Column $f$: the exact value (the conjectured value when the search below it
is complete). ``Edges searched'': the edge counts below the conjectured value that can occur. ``Graphs below'': the
candidate graphs generated at those edge counts (with the generator in use; the counts after
Lemma~\ref{lem:prune} are smaller). ``Extremal graphs'': the qualifying graphs at the conjectured value when that edge
count was also enumerated; otherwise the attaining construction \cite{BJ}.}
\label{tab:values}
\end{table}

\section{Scope}

The open window is infinite; these computations cover its members of order at most 23 and nothing beyond. The
results are exact finite statements, verified by machine; the general second line of the conjecture remains open.
The extremal graphs found at the conjectured values are listed in the computational record \cite{Record}.

\paragraph{Data availability.}
Code, records and the lists of extremal graphs are in the repository \cite{Record}, experiments EXP-005 to EXP-008
under \path{problems/combinatorics/bougard-joret/experiments/}.

\paragraph{Declaration of generative AI use.} Generative AI tools were used for literature triage, derivation
exploration, code drafting and internal review of the text. The author reviewed and verified all content and takes
full responsibility for the work.

{\small
\begin{thebibliography}{9}
\bibitem{BJ} N. Bougard and G. Joret.
\emph{Tur\'an's theorem and $k$-connected graphs}.
Journal of Graph Theory 58 (2008), 1--13.
\href{https://doi.org/10.1002/jgt.20289}{doi:10.1002/jgt.20289}.
\bibitem{DG} J. Das and S. Gupta.
\emph{Counterexample to the Bougard--Joret Conjecture}.
\href{https://arxiv.org/abs/2608.18828}{arXiv:2608.18828v1}, 2026.
\bibitem{Shell} F. Santib\'a\~nez-Leal.
\emph{The first interior shell of the Bougard--Joret extremal problem}.
CAOS Research preprint, 2026.
\href{https://doi.org/10.5281/zenodo.22315251}{doi:10.5281/zenodo.22315251}.
\bibitem{Next} F. Santib\'a\~nez-Leal.
\emph{The next matching level for triangle-free graphs with bounded degree and an exact Bougard--Joret value}.
CAOS Research preprint, 2026.
\href{https://doi.org/10.5281/zenodo.22343021}{doi:10.5281/zenodo.22343021}.
\bibitem{Turan} P. Tur\'an.
\emph{Eine Extremalaufgabe aus der Graphentheorie}.
Matematikai \'es Fizikai Lapok 48 (1941), 436--452.
\bibitem{nauty} B. D. McKay and A. Piperno.
\emph{Practical graph isomorphism, II}.
Journal of Symbolic Computation 60 (2014), 94--112.
\href{https://doi.org/10.1016/j.jsc.2013.09.003}{doi:10.1016/j.jsc.2013.09.003}.
\bibitem{Record} F. Santib\'a\~nez-Leal.
\emph{CAOS Research, Bougard--Joret EXP-005 to EXP-008: the open second regime}.
Code, enumeration records and extremal graphs, 2026.
\href{https://github.com/fsantibanezleal/CAOS_RESEARCH/tree/main/problems/combinatorics/bougard-joret/experiments}{Public experiment records}.
\end{thebibliography}
}
\end{document}
